\chapter{Introduction}
\label{ch:intro}

\section{The Problem: Contact Is Not Consequence}

Formal verification has always separated the act of writing a proof from
the act of the kernel accepting it. What has changed in the last several
years is who, or what, is doing the writing. Large language models and
autonomous coding agents can now produce syntactically plausible Lean~4
declarations, tactic scripts, and entire modules on request. This
capability is genuinely useful, and the manufacturing system described in
this thesis was itself built with substantial LLM and agent assistance.
But it raises a question that the Lean kernel alone does not settle:
\emph{what standing does generated material have before, during, and
after it is admitted?}

Kernel acceptance is a strong guarantee of one specific thing: that a term
inhabits its stated type. It is not a guarantee that the stated type is
the intended one, that a multi-part statement was not silently weakened,
that a proof obligation was not narrowed by an extra hypothesis, or that
the artifact as a whole is a faithful rendering of some target
mathematical theory. Meek et al.~\cite{meek2026formalizing} catalog
exactly these failure modes in an agent-driven numerical-analysis
pipeline: compilation success can coexist with an added weakening
hypothesis, an incompletely stated multi-part theorem, or a parameter
restriction that trivializes the intended claim. None of these failures
are kernel failures. They are failures of the surrounding process --- the
part of the pipeline that decides which candidate gets typed, checked,
and recorded in the first place.

This repository's governing position is that an LLM, an agent, a
template engine, or a human editor is, with respect to the formal
corpus, an \emph{untrusted candidate producer}. AGENTS.md states this as
the repository's Prime Directive: ``This repository is a mathematical
manufacturing system. Do not hand-code manufactured outputs. Claude is an
untrusted candidate producer. Claude may propose changes to source
declarations, templates, fixtures, tests, and hand-authored prose. Claude
may not confer standing by directly editing generated artifacts.'' The
directive is deliberately mechanical rather than aspirational: it names
which files may be hand-edited (source declaration catalogs in Turtle,
projection templates, builder scripts, stable prose) and which may not
(rendered Lean modules, ledgered paper fragments, release manifests,
release hashes, and counts). The asymmetry is the point. A candidate
producer --- human or machine --- may contact the system at many points,
but contact does not by itself create a standing claim. Only a declared,
mechanically checked admission path can do that.

Put in the vocabulary this thesis adopts throughout: the standing law is
\[
R_B \;\vdash\; A = \mu(O^*_B),
\]
read inside a declared boundary~$B$: an artifact~$A$ with standing is the
lawful manufacturing transform~$\mu$ applied to admitted observations
$O^*_B$, and the receipt~$R_B$ is what makes that consequence checkable
rather than merely asserted. Raw candidate output --- a draft theorem
statement, a proposed tactic script, a hand-typed release count --- is
$O$, not $O^*$; it acquires no standing until it passes through the
admission gate that produces $O^*_B$. This separates \emph{contact}
(a candidate producer touching the corpus) from \emph{consequence} (a
kernel-checked, receipted, replayable change of standing). The remainder
of this thesis is an account of what it takes to make that separation
operational at the scale of a several-hundred-declaration formal
library, rather than as a slogan.

\section{Thesis Statement}

This thesis argues that formal artifacts produced with the assistance of
untrusted candidate producers can be given a disciplined, replayable
notion of \emph{standing} by treating formalization itself as a
manufacturing pipeline rather than as a sequence of ad hoc edits. The
pipeline has five stages, each of which produces or checks a distinct
kind of evidence rather than merely restating the previous stage's
output:

\begin{enumerate}
\item \textbf{Source declaration.} A domain's obligation graph --- the
  structures, theorem surfaces, module dependencies, fixtures, and
  citations a formal library must cover --- is curated as an RDF/Turtle
  catalog rather than written directly as Lean source. The catalog is the
  place where scope, naming, and status (``stated'' versus a status that
  a passing guard may promote toward ``proven'') are declared.
\item \textbf{ggen projection.} The declaration catalog is projected
  through SPARQL-in-Tera templates into candidate Lean~4 modules. This
  projection step is not itself an admission step: its output is a
  candidate, carrying provenance and a content hash, with no standing of
  its own.
\item \textbf{Lean kernel admission.} The projected modules are compiled
  under a pinned toolchain. Kernel acceptance, a \texttt{sorry} census,
  and an axiom audit together decide whether a candidate becomes an
  admitted, checked declaration --- and every failure mode is a typed
  refusal rather than a silently absorbed error.
\item \textbf{mfact certification.} A release manifest aggregates the
  admitted corpus's hashes, axiom footprints, fixture results, and
  proven/stated status into a single, machine-checkable object, closed
  against a declared vocabulary of admissible objections.
\item \textbf{Replayable receipt.} The manifest, the axiom audit log, and
  the negative fixtures together constitute a receipt: evidence that a
  claimed artifact is a genuine consequence of the pipeline, checkable
  again by a third party rather than trusted on the strength of the
  first run.
\end{enumerate}

This thesis further argues that this discipline is not confined to
abstract pipeline design: it is exercised on a substantial first
manufactured vertical, \textsf{procint}, a Lean~4 library for process
intelligence --- Petri nets, workflow-net soundness, event logs and
alignment-based conformance, and OCEL-style object-centric event
data --- manufactured from the declared obligation graph rather than
hand-written declaration by declaration. Process intelligence is chosen
not because it is the terminus of the method but because, as
paper/main.tex's Scope paragraph is explicit in stating, its public
objects (events, traces, transitions, guards, conformance obligations)
already carry mathematical structure suitable for this kind of projection
and admission, making it a demonstration vertical for a method the thesis
argues is field-independent in its structure, if not in this document's
empirical coverage.

\section{Contributions}

This thesis makes the following contributions, each developed in a later
chapter:

\begin{enumerate}
\item \textbf{A typed admission discipline for formal artifacts
  (\emph{mathematical manufacturing}).} Generated output, human patches,
  templates, and agent proposals are all treated uniformly as
  candidates; standing is acquired only through an explicit, fail-closed
  admission gate --- kernel acceptance, a zero-\texttt{sorry} census, an
  axiom audit against an authorized set, negative fixtures, and manifest
  completeness (Chapter~\ref{ch:background}, Chapter~\ref{ch:mfact}).
\item \textbf{\textsf{mfact}, a Lean~4 manufacturing framework} that
  reifies candidates, refusals, admissions, manifests, certified
  releases, and objections as first-class typed values rather than as
  informal process conventions, and that discharges a per-release
  ``no valid objection'' theorem against a closed inductive family of
  admissible objection forms (Chapter~\ref{ch:mfact}).
\item \textbf{\textsf{procint}, the first manufactured vertical}: a
  process-intelligence formalization covering Petri nets, workflow-net
  soundness (including the crown-jewel soundness equivalence, proven
  under a \texttt{[Finite T]} hypothesis, and a deliberately unbounded
  countermodel witnessing why that hypothesis cannot be dropped), event
  logs, alignment-based conformance, and OCEL~2.0 object-centric logs,
  together with an account of the exactly one declaration in the current
  \TotalDecls-declaration catalog that remains genuine open proof work
  (Chapter~\ref{ch:procint}).
\item \textbf{A correspondence rail from extracted implementation to
  admitted semantics}, using Aeneas/Charon to extract a bounded Rust core
  into a proof-facing image and relating that image to admitted
  \textsf{procint} semantics under a declared, kernel-checked abstraction
  function, rather than by prose or convention alone
  (Chapter~\ref{ch:correspondence}).
\item \textbf{A manifest-generated, kernel-closed evaluation}: the
  paper's own numeric claims are rendered from the release manifest, not
  asserted in prose, and closed by a kernel-checked Standing Quadrature
  witness that refuses the release if the declaration catalog, the
  admitted corpus, and the manifest ever diverge (Chapter~\ref{ch:empirical},
  Chapter~\ref{ch:evaluation}).
\item \textbf{An account of the pipeline's own construction} by agentic
  and LLM-assisted development, examined under the same untrusted-producer
  discipline the pipeline enforces on its mathematical content
  (Chapter~\ref{ch:ai-development}).
\end{enumerate}

\section{Scope and Non-Claims}

Consistent with the Scope paragraph of the accompanying paper, this
thesis does not claim to formalize process intelligence in its entirety,
nor does it claim the manufacturing pipeline's benefits transfer to every
formal domain without further work; the generalization argument
(Chapter~\ref{ch:discussion}) is presented as a structural analogy backed
by one worked vertical, not as an empirical result over multiple
verticals. It does not claim that generated mathematics is correct by
virtue of being generated, or that kernel admission settles the harder
semantic question of whether an admitted statement is a faithful
rendering of its informal target --- a residual risk this thesis returns
to explicitly, following Meek et al.'s diagnosis, in
Chapter~\ref{ch:discussion}. It does not claim that every declaration in
the corpus is proven: the release manifest itself distinguishes
kernel-admitted, axiom-audited theorems from declarations whose formal
statement exists but whose proof obligation remains open, and this
thesis is careful to preserve that distinction rather than to elide it in
narrative prose. Where a number is asserted --- a declaration count, a
theorem count, a hash --- it is drawn from a rendered, ledgered fragment
of the release manifest rather than typed by hand into this document.

\section{Roadmap}

The remaining chapters develop the contributions above in the following
order. Chapter~\ref{ch:background} situates this work against the Lean~4
ecosystem, Petri-net and workflow-net theory, the process-mining
literature, LLM-assisted theorem proving, and Rust verification and
extraction, closing with an explicit statement of the research gap this
thesis addresses. Chapter~\ref{ch:mfact} develops the \textsf{mfact}
framework itself: candidates, typed refusals, admissions with covers
proofs, manifests, certified releases, and the closed family of valid
objections. Chapter~\ref{ch:procint} presents \textsf{procint} as the
first manufactured vertical, including its module architecture, the
workflow-net soundness development, and the deliberate countermodel that
bounds the crown-jewel theorem's scope. Chapter~\ref{ch:correspondence}
develops the Aeneas/Charon correspondence rail between extracted Rust
implementation and admitted process semantics. Chapter~\ref{ch:empirical}
reports the manifest-generated evaluation and the Standing Quadrature
closure argument in detail, and Chapter~\ref{ch:ai-development} turns the
thesis's own untrusted-producer discipline onto the agentic process that
built the pipeline. Chapter~\ref{ch:evaluation} consolidates the
quantitative release evidence referenced throughout the thesis into a
single evaluation treatment, and Chapter~\ref{ch:discussion} discusses
limitations, the residual semantic-faithfulness risk, and the
generalization argument to ontology-bearing fields beyond process
intelligence. Chapter~\ref{ch:conclusion} closes the thesis.
